\section{Conclusion}
\label{sec:conclusion}
%In this work, we introduced graph-based conformal compression, a framework that bridges the statistical rigor of conformal prediction with the algorithmic toolkit of combinatorial optimization. %By formulating the compression task as a density-constrained subgraph selection problem, we showed that the problem of generating high-coverage subsets of weighted hypergraphs admits efficient bicriteria approximations via linear programming. Our experiments on synthetic navigation and trip-planning scenarios demonstrate that this approach resolves structural ambiguities, such as ``honey pot'' traps and bottlenecks, where standard frequency-based heuristics fail.
Our framework opens several avenues for future research. First, many practical domains exhibit specific hypergraph structures. For instance, road networks are nearly planar, and hierarchical classification labels form trees. It is an open question whether the conformal subgraph problem admits tighter approximations (e.g., PTAS) or exact solutions when restricted to these special graph classes. Similarly, real-world predictive distributions often exhibit  planted structure, like the trip planning scenario in Appendix~\ref{app:trip}. Can we prove tighter bounds or exact recovery guarantees for our algorithms under these generative assumptions?

Stepping back, we are treating the predictive model as fixed. However, conformal validity holds for \textit{any} scoring function. This allows for an end-to-end approach where the model $f_\theta$ is trained to explicitly optimize for compressibility, One could, for instance, augment the training objective with a differentiable ``graph compression'' regularizer. However, such regularizers will introduce significant non-convexity, addressing which forms an interesting research direction.

%Finally, even when treating the predictive model as fixed,  we have decoupled the conformal guarantee from the compression guarantee, which typically holds only if the predictive model approximates the true posterior. A compelling avenue for future work is to directly optimize the conformal set for efficiency, rather than calibrating a fixed model. This is analogous to~\cite{gao2025volume} who provide finite-sample guarantees on the volume of the conformal set, albeit requiring the calibration dataset to be large enough to learn the geometry of the optimal set. Applying such frameworks to our setting, where the calibration data is highly sparse and high-dimensional, remains a challenging open problem.